\documentclass[a4paper]{article}     
\usepackage{amsfonts}
\usepackage{amsmath}
\usepackage{amssymb}
\usepackage{graphicx}


\begin{document}

\noindent \large Auteur: Adsayan Gunaseelan \\
\noindent \large Studierichting: Informatica\\
\noindent \large Oefeningenreeks 2E nr. 10.15\\

\medskip

\normalsize
10.15) Los de volgende exponentiële vergelijkingen op door gebruik van een gepaste substitutie:\\

$1+3^{2x-3} = 4\cdot3^{x-2}$\\

$\Leftrightarrow 1+3^{2(x-2)+1} = 4\cdot3^{x-2}$\\

$\Leftrightarrow 1+3\cdot3^{2(x-2)} = 4\cdot3^{x-2}$\\

$\Leftrightarrow 1+3y^2 \overset{y=3^{x-2}}{=} 4y$\\

$\Leftrightarrow 3y^2 -4y +1 = 0$\\

$D = 16 - 4\cdot 3 \cdot 1 = 16 - 12 = 4, \sqrt{D} = 2$\\

$y_1 = \dfrac{4-2}{6} = \dfrac{1}{3}, y_2 = \dfrac{4+2}{6} = 1$\\

$\Rightarrow y_1$:\\

$3^{x-2} = \dfrac{1}{3}$\\

$\Leftrightarrow 3^{x-2} = 3^{-1}$\\

$\Leftrightarrow x-2 = -1$\\

$\Leftrightarrow x = 1$\\


$\Rightarrow y_2$:\\

$3^{x-2} = 1$\\

$\Leftrightarrow 3^{x-2} = 3^{0}$\\

$\Leftrightarrow x-2 = 0$\\

$\Leftrightarrow x = 2$\\

\begin{thebibliography}{999}
%\bibitem{a} Referentie
\end{thebibliography}
\end{document} 